% Generated draft: requires mathematical and source review.
\documentclass[11pt]{article}
\usepackage[margin=25mm]{geometry}
\usepackage{fontspec}
\setmainfont{Latin Modern Roman}
\usepackage{amsmath,amssymb,mathtools,unicode-math}
\setmathfont{Latin Modern Math}
\usepackage{enumitem,needspace,xurl,hyperref}
\setlength{\parindent}{0pt}
\setlength{\parskip}{5pt}
\setlength{\emergencystretch}{3em}
\newcommand{\N}{\mathbb{N}}
\newcommand{\Z}{\mathbb{Z}}
\newcommand{\Q}{\mathbb{Q}}
\newcommand{\R}{\mathbb{R}}
\newcommand{\C}{\mathbb{C}}
\newcommand{\F}{\mathbb{F}}
\newcommand{\A}{\mathbb{A}}
\newcommand{\PP}{\mathbb{P}}
\newcommand{\E}{\mathbb{E}}
\newcommand{\eps}{\varepsilon}
\newcommand{\dd}{\,\mathrm{d}}
\begin{document}
\section*{UnsolvedMath ID: 30003054}
\textbf{Source problem code:} catalogue entry
\subsection{Definitions and mathematical statement}
A Koopman operator acts on observables by composition with a dynamics map: for a map \(T:X\to X\) and a scalar-valued observable \(f\), its action is \(Uf=f\circ T\), whenever this defines an operator on the chosen observable space. An eigenfunction of \(U\) is a nonzero observable \(\phi\) satisfying \(U\phi=\lambda\phi\) for some scalar \(\lambda\). This notation explains the terminology only; the supplied record does not specify the dynamics, the scalar field, or whether the construction concerns a single operator, a family of operators, or a generator.

Write \(A_\alpha f\) schematically for the generalized Laplace averages referenced in the source, with \(\alpha\) denoting their unspecified limiting parameter. The record reports a construction of Koopman eigenfunctions using convergence in an unspecified observable-space norm. If \(B\) denotes that space and \(\phi\) the asserted limit, the reported convergence has the schematic form
\[
\lVert A_\alpha f-\phi\rVert_B\longrightarrow 0.
\]
\textbf{Question 1.} Under the hypotheses of that construction, does the same family converge pointwise to the eigenfunction, in the sense intended by the source,
\[
A_\alpha f(x)\longrightarrow\phi(x)?
\]
The supplied record does not specify whether this is required at every state, almost everywhere with respect to a particular measure, or on a distinguished subset. If observables are equivalence classes, the representatives needed for pointwise evaluation are also unspecified.

\textbf{Question 2.} For which other spaces of observables does the same generalized Laplace-average construction yield Koopman eigenfunctions by norm convergence, especially when the norms are readily computable?

These are source-dependent questions rather than a fully specified quantified conjecture. The averaging formula, limiting regime, original convergence theorem and its hypotheses, admissible observables, spectral parameters, and criteria for a nonzero eigenfunction limit are absent from the supplied record. Consequently, neither the scope of the pointwise assertion nor the precise property required of alternative observable spaces can be fixed without additional source information.
\subsection{Short English statement}
The problem asks whether the generalized Laplace averages in a reported norm-convergent construction of Koopman eigenfunctions also converge pointwise, and which other observable spaces support that construction with readily computable norms. The supplied record omits the construction and its hypotheses, so both questions require further specification before they become precise conjectures.
\subsection{Sources}
\begin{itemize}
\item Supplied local UnsolvedMath registry record, permanent ID \(30003054\), title \emph{Pointwise Convergence of Koopman Laplace Averages}; exact source problem code \texttt{catalogue entry}. Collection: Oberwolfach Reports Open Problems; proposed year: 2016. No individual proposer is identified.
\item Supplied report citation: \emph{Mini-Workshop: Applied Koopmanism} (2016), Oberwolfach Reports, pp. 297--340, DOI 10.4171/owr/2016/7. Source URL and supplied literature-source link labelled ``original OWR report'': \url{https://doi.org/10.4171/owr/2016/7}.
\item Literature assessment dated exactly \texttt{2026-08-21}: the supplied curation labels the problem open, reports that the source explicitly leaves pointwise convergence open, and asks about other observable spaces supporting usable norm convergence. This is supplied dated curation, not an independently verified literature assessment. The record itself requests independent verification of time-sensitive status and cited claims.
\item The supplied statement status is \texttt{corrected\_verified}, accompanied by a claim of checking against the official report and preserving its hypotheses, notation, and target. That verification claim is part of the supplied provenance; the report has not been independently consulted for this draft.
\end{itemize}
\end{document}
